
\begin{algorithm}[!t]\caption{\small The nearest-neighbor chain algorithm for HAC.}\label{alg:nnchain}
\KwIn{$G = (V, E, w)$}

\SetKwFunction{FNNChain}{NearestNeighborChain}
\SetKwProg{Fn}{Function}{:}{}

\Fn{\FNNChain{$G$}}{
$\clustering \gets$ clustering where each $v \in V$ is in a separate cluster and is {\em active}.\\
\For{each cluster $v \in V$}{
  \If {$v$ is active}{
    Initialize a stack $S$ initially containing only $v$ \\
    \While {$S$ is not empty}{
      $t \gets \textsc{Top}(S)$.\\
      $(t, b, w_{t,b}) \gets \textsc{BestEdge}(t)$.\label{line:nnchainbestedge} \\
      \If{$b$ is already on $S$}{
        $\textsc{Pop}(S)$. \\
        $\textsc{Merge}(t, \textsc{Top}(S))$.\label{line:nnchainmerge}  \ \  //\ {\em $t$ is marked as inactive; $\textsc{Top}(S)$ remains active.}\\
        $\textsc{Pop}(S)$.
       } \Else {
        $\textsc{Push}(S, b)$.
       }
    }
  }
}
}
\end{algorithm}

\section{\texorpdfstring{\boldmath$O(mh\log n)$}{O(mhlogn)} Upper Bound}\label{sec:workupperbound}

%\knote{I think we need to define 'active' clusters. Also, I think the pseudocodes never mark clusters as 'inactive'.}

In this section we show that we can solve graph-based average linkage HAC in $\tilde{O}(mh)$ time where $h$ is the {\em height} of the output dendrogram.
Thus, when $mh = o(n^{3/2})$ this algorithm improves over the lower bound from Section~\ref{sec:worklb}.
Interestingly, we give a simple analysis showing that either of the two classic approaches to HAC---the nearest-neighbor chain or heap-based algorithms---achieves this bound.
We give the pseudocode for the nearest-neighbor chain algorithm in Algorithm~\ref{alg:nnchain} and the heap-based algorithm in Algorithm~\ref{alg:heapbased}.
In a nutshell, in both algorithms initially all vertices start in their own cluster, and are {\emph{active}}.
As the algorithms proceed, they find edges to merge that the exact greedy HAC algorithm (Algorithm~\ref{alg:staticgraph}) will make and merge them, marking one of the endpoints as {\em inactive}.
This process continues until no further edges between active clusters remain.
We discuss the algorithms individually in more detail in what follows.

We start by discussing some data structures and details in how the algorithms carry out merges.

\paragraph{Neighborhood representation}
The neighbors of each vertex are stored in a max-heap supporting standard operations (e.g., \textsc{find-min}, \textsc{delete-min}, \textsc{insert}, \textsc{decrease-key}).
The priority of a neighbor is the average linkage weight of the edge to the neighbor.
At the start of the algorithm, the heaps are initialized with the initial neighbors of each vertex.
Let $N(A)$ represent the neighbors of a cluster $A$.
For concreteness in what follows, we use Fibonacci heaps~\cite{CLRS, fredman1987fibonacci}.


\paragraph{Merging clusters}
Next, we discuss how to implement the {\em merge} step in both algorithms, which merges two clusters $A$ and $B$, and their corresponding heaps (i.e., merging $N(A)$ and $N(B)$ and updating the average linkage weights of {\em all} edges affected by the merge).
We maintain the invariant that the heaps always store the correct average linkage weights of all edges.
We first compute $\mathcal{I}(A,B) = N(A) \cap N(B)$, which can be done by sorting both sets and merging the sorted arrays.
For each $C \in \mathcal{I}(A, B)$, without loss of generality, we remove the reference to $A$ in the heap and update the priority of $B$ to reflect the new average linkage weight from $C$ to $A \cup B$ using \textsc{decrease-key}.
Similarly, we update each $C \in N(A) \setminus \mathcal{I}(A, B)$ to reference $B$ instead of $A$.
Finally, for each $C \in N(A) \cup N(B)$ we update the priority of the edge $(C, B)$ to reflect the new size of $B$.

At the end of the merge: (1) $N(B)$ contains $N(A) \cup N(B) \setminus \{A, B\}$; (2) all edges in the updated $N(B)$ have their average linkage weight correctly set; and (3) after the update all $C \in N(B)$ no longer reference $A$, but point to $B$ and have their average linkage weights correctly set.
%\hnote{Shall we explicitly say that the two heaps are merged into one?}
%
The cost of this merge procedure is dominated by the cost of sorting $N(A)$ and $N(B)$; plugging in the cost bounds for any reasonable heap data structure (e.g., a Fibonacci heap) yields the following lemma:
\begin{lemma}\label{lem:merge}
Merging two clusters $A,B$ with $N(A) + N(B) = T$ can be done in $O(T \log n)$ time.
\end{lemma}


\paragraph{Nearest-Neighbor Chain Analysis}
As we run the nearest-neighbor chain algorithm (Algorithm~\ref{alg:nnchain}), the algorithm only requires two non-trivial operations beyond basic data structures (e.g., stacks and heaps):\footnote{In fact the same two primitives are also the only non-trivial ones used for the heap-based algorithm (Algorithm~\ref{alg:heapbased}) and thus we fully specify the details for this algorithm as well.}
%\knote{We need to add the algorithm pseudocode for NN chain here}:
\begin{enumerate}
  \item \textsc{BestEdge}($C$): fetch the highest similarity edge incident to a cluster $C$.
  \item \textsc{Merge}($A, B$): merge the clusters $A,B$ into a cluster $A \cup B$ (represented, without loss of generality by $B$). Thus after the merge, $A$ is inactive, and $B$ remains active.
%  \knote{I think this means that $A$ becomes inactive, not, as we say in pseudocodes, one of A, B, right?}
\end{enumerate}

Using the data structures above, we can implement both operations very quickly. 
In particular, Lemma~\ref{lem:merge} gives the time for merging two clusters, and \textsc{BestEdge}($C$) can be implemented in $O(1)$ time using the \textsc{find-min} operation.

After merging two clusters that are reciprocal ``best'' neighbors, the algorithm back-tracks to the third vertex in the chain and restarts the process, using calls to \textsc{BestEdge}.
The cost of these calls can be charged to the merges, resulting in a total of $O(n)$ time for the $n-1$ merges.
The time that still must be accounted for is that of \textsc{Merge}.

\begin{claim}\label{claim:mergecost}
The total amount time spent to perform all merges made by the nearest-neighbor chain algorithm is $O(mh\log n)$.
\end{claim}
\begin{proof}
Split each edge $(u,v)$ into two directed edges $(u,v)$ and $(v,u)$, since the edge will be present initially in both $N(u)$ and $N(v)$ (and the intermediary clusters they form until they are merged).

Consider a directed edge $(u,v)$, and the merges it experiences as it flows along the path from the initial cluster containing it ($u$) to the cluster where it is finally merged.
The dendrogram that the algorithm constructs has height $h$ by definition, and thus the length of this path is at most $h$.
In each node on the path, the edge contributes a cost of $O(\log n)$ to the merge cost in Lemma~\ref{lem:merge}.
Thus, the cost of this edge across the entire path is at most $O(h \log n)$.

Applying this charging argument for all $2m$ directed edges, the total cost over all merges is $O(mh\log n)$.
\end{proof}


Putting together the above yields the following theorem.
\mainUpperSec*


\begin{algorithm}[!t]\caption{The heap-based algorithm for HAC.}\label{alg:heapbased}
\KwIn{$G = (V, E, w)$}

\SetKwFunction{FHeapBased}{HeapBasedHAC}
\SetKwProg{Fn}{Function}{:}{}

\Fn{\FHeapBased{$G$}}{
$\clustering \gets$ clustering where each $v \in V$ is in a separate cluster and is {\em active}.\\
Let $H$ be a max-heap storing the highest-weight edge incident to each active cluster in the graph \\
\While{$|H| > 1$}{
  Let $e=(u, v, w_{u,v})$ be the max weight edge in $H$. \\
  Delete $e$ from $H$. \\
  \If {$v$ is inactive} { \label{line:heapbadextract}
    Let $e' = (u, v', w_{u,v'}) = \bestedge{u}$. \\
    Insert $e'$ into $H$. \label{line:heapinsertfirst} \\
  } \Else {
    $x = \textsc{Merge}(u, v)$.    \ \  //\ {\em $u$ is marked as inactive; $v$ remains active.} \\
    Let $e' = (x, y, w_{x,y}) = \bestedge{x}$. \\
    Insert $e'$ into $H$. \\
  }
}
}
\end{algorithm}

\paragraph{Heap-Based Algorithm}

We note it is easy to extend the argument above to obtain a similar bound for the heap-based algorithm 
(pseudocode shown in Algorithm~\ref{alg:heapbased}).
In particular, beyond the cost for merging, which is the same in the nearest-neighbor chain algorithm and the heap-based algorithm (captured by Claim~\ref{claim:mergecost}), the only extra work incurred by the heap-based algorithm is due to edges unsuccessfully extracted from the heap (Line~\ref{line:heapbadextract} of Algorithm~\ref{alg:heapbased}). 
The number of such unsuccessful edges is at most $m$, and each such edge incurs $O(\log n)$ work due to an extra call to \bestedge{u}, for a total of $O(m\log n)$ extra work.
Thus the total work of the heap-based algorithm is also $O(mh\log n)$.

\begin{theorem}
There is an implementation of the heap-based algorithm for average linkage HAC that runs in $O(mh \log n)$ time where $h$ is the height of the output dendrogram.
\end{theorem}


%\knote{Technically, Alg 1 does not use a heap.}




